\section{Importancia de la caracterización geotécnica del suelo}

La respuesta de una obra civil frente a cargas estáticas y dinámicas depende de propiedades mecánicas del terreno que varían con la profundidad y, en general, también lateralmente. Reconocer esa variación (espesores, contactos, rellenos y zonas de menor rigidez) es una condición previa al dimensionamiento de fundaciones, a la evaluación de asentamientos y a la estimación de la respuesta sísmica de sitio.

Entre las magnitudes que describen ese comportamiento, la velocidad de propagación de las ondas de corte, $V_S$, ocupa un lugar particular porque se vincula de forma directa con la rigidez del medio a pequeñas deformaciones. Para un medio elástico de densidad $\rho$, el módulo de corte a pequeñas deformaciones queda determinado por

\begin{equation}
G_{\max}=\rho\,V_S^{2},
\end{equation}

de manera que medir $V_S$ equivale, conocida la densidad, a medir la rigidez del suelo en el rango de deformaciones en que su comportamiento puede considerarse aproximadamente elástico. Esa equivalencia explica por qué el perfil $V_S(z)$ se ha consolidado como parámetro de entrada en la clasificación sísmica de sitios y en el análisis de respuesta dinámica del terreno \citep{Kramer1996,Foti2014}.

Conviene fijar desde el inicio una precisión que condiciona todo lo que sigue: ninguna técnica de prospección entrega una imagen exacta del subsuelo, sino un modelo compatible con los datos, con la geometría del ensayo y con los supuestos de interpretación. Esa distinción determina qué información debe preservar el instrumento y hasta qué profundidad puede sostenerse una afirmación.

Se tiene como objetivo de este proyecto de fin de grado el diseño, la implementación y la evaluación experimental de un sistema electrónico multicanal del tipo IoT para la caracterización geotécnica del suelo mediante ondas sísmicas, orientado a profundidades de investigación de hasta 50 m. La contribución que se presenta en esta primera etapa es de instrumentación: la integración y la evaluación con datos reales de una cadena de medida completa, desde la fuente instrumentada y el geófono hasta el acondicionamiento analógico, la conversión analógico-digital, la temporización, el almacenamiento y el procesamiento que conduce a un perfil $V_S(z)$ preliminar. El interés de construirla en lugar de adquirirla no está en suponer que la industria carezca de soluciones, sino en disponer de una plataforma auditable en cada etapa, condición que el instrumental comercial cerrado no ofrece a un laboratorio universitario.

Conviene precisar el centro de gravedad del trabajo antes de continuar. Se trata de un proyecto de fin de grado de Ingeniería Electrónica, de modo que el objeto de diseño, implementación, caracterización y validación es el sistema de instrumentación. La caracterización geotécnica y el método MASW constituyen el contexto de aplicación: son los que definen los requerimientos que el instrumento debe satisfacer y el criterio con el que se juzga si los satisface. Las secciones 2 a 7 desarrollan en consecuencia únicamente la física necesaria para derivar esos requerimientos, y la interpretación geológica se mantiene en todo el documento dentro del límite estricto de verificar la plausibilidad de los resultados obtenidos.

Ese objetivo formal debe distinguirse con claridad del alcance efectivamente alcanzado hasta hoy, y el documento mantiene esa separación en todas sus secciones:

\begin{itemize}
\item \textbf{Objetivo de diseño.} Una plataforma multicanal orientada a investigar hasta 50 m de profundidad.
\item \textbf{Alcance demostrado con la campaña y la cadena actuales.} Adquisición, procesamiento e inversión con soporte de investigación del orden de 10 a 11 m.
\item \textbf{Brecha experimental.} Aumentar la energía coherente de baja frecuencia, completar el hardware multinodo y realizar las validaciones metrológicas y geofísicas independientes que todavía faltan.
\end{itemize}

Las secciones 2 a 7 establecen qué se desea medir y por qué mediante ondas superficiales; la sección 8 traduce esa elección en requerimientos cuantitativos; la sección 9 presenta el diseño electrónico como respuesta a ellos; las secciones 10 y 11 describen la validación y sus resultados, y la 12 delimita lo que aún no puede afirmarse.

\section{Métodos tradicionales de caracterización del suelo}

El reconocimiento geotécnico convencional se apoya en ensayos que establecen contacto físico con el terreno. El ensayo de penetración estándar (SPT) da una medida indirecta de resistencia y permite recuperar muestras a intervalos discretos, con la ventaja de una práctica extendida y correlaciones documentadas; el ensayo de penetración con cono (CPT y su variante piezométrica CPTu) mejora la continuidad vertical del registro y reduce la dependencia del operador, aunque pierde aplicabilidad en gravas o niveles cementados.

Cuando el objetivo es directamente una velocidad sísmica, la referencia son los ensayos en pozo: el downhole registra en profundidad las llegadas generadas en superficie y el crosshole mide el tiempo de tránsito entre perforaciones vecinas, ofreciendo la mejor resolución vertical disponible para $V_S$. Ambos comparten la limitación que motiva este trabajo, ya que exigen perforar, con el costo y el plazo asociados, y describen el suelo sólo en el entorno del pozo.

De esa comparación se desprende la motivación de los métodos no invasivos. Una obra necesita reconocer la variación lateral de la rigidez en un volumen extenso, y un conjunto de sondeos puntuales resulta costoso de densificar. Los métodos basados en la propagación de ondas mecánicas desde la superficie permiten cubrir esa extensión sin perforar, a cambio de trabajar con una magnitud que se infiere en lugar de medirse por contacto. A ese argumento se agrega un campo de aplicación que el carácter no invasivo habilita de manera casi exclusiva: el control de calidad y el seguimiento de obras ya construidas. Una traza vial, un terraplén o una plataforma pueden recorrerse de forma repetida y sin daño para verificar la uniformidad de la compactación durante la construcción, y para detectar posteriormente la degradación o la pérdida de rigidez que anticipa la necesidad de mantenimiento. Un ensayo destructivo no admite esa repetición sobre la misma obra.

\section{Métodos basados en ondas mecánicas}

La caracterización sísmica se apoya en que las propiedades mecánicas del medio gobiernan cómo se propaga una perturbación, y en que esa propagación deja magnitudes observables en superficie: tiempos de llegada, amplitudes y relaciones de fase entre puntos separados. Medirlas permite estimar las propiedades sin acceder a la profundidad de interés.

Las familias de métodos se distinguen por qué observable explotan. La refracción sísmica utiliza los tiempos de primera llegada de ondas de cuerpo refractadas críticamente y resulta apropiada cuando la velocidad crece de manera monótona con la profundidad; su punto débil conocido es la incapacidad de detectar capas de baja velocidad ocultas bajo otras más rápidas. La reflexión aprovecha los contrastes de impedancia y alcanza mayor profundidad, pero requiere fuentes de energía considerable y una relación señal a ruido difícil de obtener en el rango somero. Los métodos de ondas superficiales, en cambio, explotan la dependencia de la velocidad de fase con la frecuencia, y esa dependencia es precisamente lo que codifica la variación de rigidez con la profundidad.

Para un ensayo desde superficie orientado a $V_S$ en los primeros metros, la última familia presenta una ventaja física decisiva. Un impacto vertical aplicado en la superficie de un semiespacio destina la mayor parte de la energía elástica radiada a la onda Rayleigh \citep{Foti2014}, cuya amplitud decae con la distancia de forma considerablemente más lenta que la de las ondas de cuerpo. El método aprovecha así la componente más energética del campo de ondas en lugar de combatirla como ruido, lo que resulta determinante cuando la fuente disponible es un impacto manual y no una fuente sísmica de gran porte.

\section{Fundamentos de propagación de ondas en un medio}

En un medio elástico, homogéneo e isótropo, la propagación de perturbaciones infinitesimales queda descrita por dos velocidades independientes, asociadas a los dos modos posibles de deformación. La velocidad de las ondas de compresión y la de las ondas de corte se expresan en función de los parámetros de Lamé $\lambda$ y $\mu$ y de la densidad $\rho$ como

\begin{equation}
V_P=\sqrt{\frac{\lambda+2\mu}{\rho}},
\qquad
V_S=\sqrt{\frac{\mu}{\rho}} .
\end{equation}

La segunda expresión es la que conecta la medición con el objetivo planteado en la sección 1: como $\mu$ coincide con el módulo de corte a pequeñas deformaciones, obtener $V_S$ equivale a obtener $G_{\max}$ salvo por el factor de densidad. Esa relación resulta especialmente conveniente porque las dos magnitudes involucradas tienen rangos de variación muy dispares. La densidad de los suelos habituales se mantiene aproximadamente entre 1,4 y 2,2 t/m³, es decir dentro de un factor menor que dos, mientras que $V_S$ recorre desde algunas decenas de metros por segundo en suelos blandos hasta varios cientos en materiales densos o cementados. Como $G_{\max}$ depende del cuadrado de la velocidad, su variación queda dominada por $V_S$, y un valor de densidad estimado con precisión moderada resulta suficiente para obtener una rigidez útil \citep{Kramer1996,Foti2014}. Las ondas de compresión, en cambio, resultan poco informativas en el rango somero saturado: por debajo del nivel freático $V_P$ queda gobernada por la compresibilidad del agua intersticial y tiende hacia unos 1500 m/s con independencia del estado del esqueleto sólido, de modo que pierde sensibilidad justamente a la propiedad que interesa medir \citep{Foti2014,Foti2018}.

Ambas velocidades quedan vinculadas por el coeficiente de Poisson,

\begin{equation}
\nu=\frac{V_P^{2}-2V_S^{2}}{2\left(V_P^{2}-V_S^{2}\right)},
\end{equation}

relación que conviene retener porque $\nu$ no fue medido en este trabajo y debió adoptarse por hipótesis en la inversión, con las consecuencias que se analizan en la sección 11.3.

Las expresiones anteriores suponen un medio homogéneo y perfectamente elástico, y conviene señalar que ninguna de las dos condiciones se cumple en el terreno real. El suelo es un medio estratificado, y esa estratificación no es una desviación molesta del modelo sino precisamente lo que el ensayo se propone reconstruir. El suelo tampoco es perfectamente elástico: disipa energía a medida que la onda avanza, por lo que la amplitud registrada en un receptor lejano se reduce tanto por divergencia geométrica como por atenuación material, y esta última crece con la frecuencia.

De ahí se sigue una consecuencia que condiciona el diseño del instrumento. Si la señal útil se debilita al aumentar la distancia y la frecuencia, entonces el margen entre señal y ruido de la cadena de medida deja de ser una cifra de mérito genérica y pasa a ser la variable que determina cuál es la máxima longitud de onda que todavía puede reconocerse y, por lo tanto, hasta qué profundidad puede investigarse.

\section{Ondas de cuerpo y ondas superficiales}

Las ondas de cuerpo se propagan por el volumen del medio y se distinguen por su polarización: la P desplaza las partículas en la dirección de propagación y la S lo hace transversalmente, razón por la cual es esta última la que involucra al módulo de corte. Dentro de la onda S conviene además distinguir dos polarizaciones, porque de ellas depende qué onda superficial puede formarse: la componente SV oscila en el plano vertical que contiene la dirección de propagación y la SH en el plano horizontal perpendicular a él.

Las ondas superficiales no son un tercer tipo de onda independiente, sino el resultado de la interacción de las anteriores con la superficie libre del terreno. Cuando una onda P y una onda SV inciden sobre esa frontera, la condición de tensión nula obliga a que se conviertan mutuamente al reflejarse, y para ciertas combinaciones de ángulo y frecuencia ambas quedan acopladas en una perturbación que no se propaga hacia el interior sino que viaja paralela a la superficie, con amplitud que decae exponencialmente con la profundidad. Esa perturbación es la onda Rayleigh: combina por construcción una componente vertical y una longitudinal, y describe en el semiespacio homogéneo un movimiento elíptico retrógrado. Su velocidad de fase se sitúa entre el 87 \% y el 96 \% de $V_S$ según el coeficiente de Poisson, lo que la convierte en un buen sustituto observable de la velocidad de corte.

La componente SH no puede acoplarse a la P por razones de polarización, y por eso no participa de la onda Rayleigh. Puede, en cambio, quedar atrapada por reflexiones múltiples dentro de una capa superficial de menor velocidad apoyada sobre un sustrato más rápido, y forma entonces la onda Love. Como su movimiento es horizontal y transversal, no produce señal apreciable en un geófono vertical; por esa razón la onda Love queda fuera del alcance de este trabajo, cuya cadena de medida registra únicamente la componente vertical.

La relación entre la Rayleigh y $V_S$ es la que convierte a la primera en un observable útil. Como la velocidad de fase Rayleigh es del orden del 90 \% de la velocidad de corte del material que la onda efectivamente muestrea, una medición de $c_R$ constituye una medición aproximada de $V_S$ una vez resuelto qué volumen del terreno correspondió a cada frecuencia.

Esa correspondencia entre frecuencia y volumen muestreado la fija la longitud de onda. Para una frecuencia $f$ y una velocidad de fase $c_R(f)$,

\begin{equation}
\lambda_R(f)=\frac{c_R(f)}{f},
\end{equation}

y el desplazamiento asociado a una onda Rayleigh decae con la profundidad en una escala comparable a $\lambda_R$. Las longitudes de onda cortas quedan por tanto controladas por el material somero, mientras que las largas integran un volumen mayor y alcanzan mayor profundidad. La regla práctica $z\approx\lambda/2$, en la que $z$ es la profundidad de investigación, entendida como la profundidad máxima hasta la cual la medición conserva sensibilidad apreciable, resume esa proporcionalidad y se emplea de forma habitual para estimar la profundidad de investigación de una campaña, pero no asigna una profundidad exacta a cada frecuencia: las funciones de sensibilidad son distribuidas y se superponen entre frecuencias vecinas. De ahí se sigue una advertencia que resulta decisiva al interpretar los registros: observar energía a una frecuencia baja no demuestra que exista allí una onda Rayleigh útil, criterio que la sección 12 desarrolla al discutir la fuente.

\section{Utilización de ondas superficiales para la caracterización de suelos}

En un medio homogéneo la velocidad Rayleigh no depende de la frecuencia. La dispersión aparece cuando el medio está estratificado, porque cada longitud de onda muestrea una combinación distinta de capas y se propaga con la velocidad de fase que resulta de esa combinación. La curva $c_R(f)$ deja entonces de ser una constante del material y pasa a contener, codificada, la variación de rigidez con la profundidad: ese es el fundamento del método.

La explotación de esa información sigue la cadena de la figura 1: se registra la respuesta del terreno en varias posiciones de una línea, se transforma el conjunto de trazas al dominio frecuencia–velocidad de fase, se extrae de esa imagen la curva de dispersión observada $c_R^{\mathrm{obs}}(f)$ y se resuelve el problema inverso buscando el modelo estratificado que la reproduzca \citep{Park1998,Park1999,Xia1999}.

La última etapa impone la cautela principal del método: el problema inverso no tiene solución única, de modo que un ajuste bajo demuestra consistencia entre el dato y el modelo propuesto, no que ese modelo sea el único compatible. La sección 11.3 aplica ese criterio al resultado obtenido.

\begin{figure}[!tb]
\centering
\footnotesize
\begin{tikzpicture}[
  node distance=3.2mm,
  caja/.style={draw,rounded corners=1.5pt,align=center,inner sep=3pt,
               text width=0.80\columnwidth,font=\footnotesize},
  fl/.style={-{Latex[length=1.6mm]},thick}]
  \node[caja] (a) {Propiedad mecánica del suelo\\[1pt]\scriptsize $V_S(z)$, la incógnita};
  \node[caja,below=of a] (b) {Propagación de ondas Rayleigh\\[1pt]\scriptsize cada $\lambda$ muestrea un volumen distinto};
  \node[caja,below=of b] (c) {Registros multicanal $u(x,t)$\\[1pt]\scriptsize lo que el instrumento debe preservar};
  \node[caja,below=of c] (d) {Imagen frecuencia--velocidad de fase\\[1pt]\scriptsize los modos aparecen como máximos};
  \node[caja,below=of d] (e) {Curva de dispersión $c_R^{\mathrm{obs}}(f)$};
  \node[caja,below=of e] (f) {Inversión\\[1pt]\scriptsize modelo estratificado que reproduce la curva};
  \node[caja,below=of f] (g) {Perfil $V_S(z)$ con banda de incertidumbre};
  \draw[fl] (a)--(b); \draw[fl] (b)--(c); \draw[fl] (c)--(d);
  \draw[fl] (d)--(e); \draw[fl] (e)--(f); \draw[fl] (f)--(g);
  \node[right=1mm of b.east,font=\scriptsize,anchor=west,xshift=-1mm] {};
\end{tikzpicture}
\caption{Cadena de inferencia del método. Las dos relaciones que la gobiernan son
$\lambda_R=c_R/f$, que liga cada frecuencia con el volumen muestreado, y
$z\approx\lambda/2$, que traduce esa longitud de onda en profundidad de investigación.
La flecha descendente es el sentido de la medición; la inversión recorre el camino
inverso y por eso no tiene solución única.}
\end{figure}

\section{Revisión y selección del método}

\begin{table*}[!tb]
\centering
\footnotesize
\caption{Comparación de los métodos considerados. La selección de MASW activo se sigue
de la propiedad buscada y de los recursos disponibles, no de una jerarquía absoluta
entre técnicas.}
\begin{tabularx}{0.94\textwidth}{@{}lYYY@{}}
\toprule
Método & Propiedad observada & Requisito de campo & Limitación dominante \\
\midrule
SPT & Resistencia a la penetración & Perforación & Medida puntual y por correlación empírica \\
CPT / CPTu & Resistencia de punta y fricción lateral & Penetración continua & Pierde aplicabilidad en gravas y cementados \\
Downhole / Crosshole & $V_P$ y $V_S$ directas & Uno o varios pozos & Costo y plazo; describe sólo el entorno del pozo \\
Refracción sísmica & Tiempos de primera llegada & Fuente y línea en superficie & No detecta capas de baja velocidad ocultas \\
SASW & Diferencia de fase entre dos receptores & Dos receptores y ensayos repetidos & Sin criterio interno de separación modal \\
MASW activo & Dispersión $c_R(f)$ sobre varios offsets & Fuente controlada y arreglo de receptores & Depende de la energía disponible en baja frecuencia \\
Pasivos (ReMi, SPAC, ESAC) & Dispersión del ruido ambiental & Registro simultáneo en todo el arreglo & Exige simultaneidad y un campo de ruido adecuado \\
\bottomrule
\end{tabularx}
\end{table*}

Dentro de la familia de ondas superficiales, la alternativa histórica a MASW es el análisis espectral de ondas superficiales (SASW), que estima la velocidad de fase a partir de la diferencia de fase entre dos receptores y repite la medición con distintas separaciones. Su exigencia instrumental es menor, pero la estimación depende del desenvolvimiento de fase y no dispone de un criterio interno para separar el modo fundamental de los modos superiores, lo que obliga a un procesamiento cuidadoso y a la repetición del ensayo. El registro simultáneo de muchos receptores que propone MASW resuelve precisamente esa debilidad: la transformación al dominio $f$–$c$ separa los modos como máximos distintos y proporciona redundancia frente a ruido local y frente a la falla de un canal individual \citep{Park1999}.

Los métodos pasivos, entre ellos ReMi, SPAC y ESAC, presentan una ventaja complementaria de peso, ya que el ruido ambiental contiene energía en frecuencias inferiores a las que alcanza un impacto manual y permite en principio extender el rango hacia mayores profundidades. Su requisito, sin embargo, resulta incompatible con los recursos disponibles: exigen registrar simultáneamente en todas las posiciones del arreglo, porque el campo de ruido no puede reproducirse a voluntad mientras se desplaza un sensor. Con dos geófonos disponibles en la Facultad y un único nodo receptor operativo por registro, estos métodos quedan fuera de alcance por construcción y no por decisión de diseño.

La selección de MASW activo se sigue entonces de los objetivos y de la física expuesta. La fuente controlada permite referenciar temporalmente cada registro y repetir la excitación tantas veces como sea necesario; la geometría conocida convierte la diferencia de fase entre posiciones en una velocidad; y el registro de múltiples offsets ofrece la redundancia que la inversión necesita. Para el propósito de este trabajo hay además una razón instrumental: MASW ejercita simultáneamente todas las funciones que el sistema debe cumplir, esto es, sincronización, preservación de fase entre canales, margen dinámico, almacenamiento sin pérdidas y trazabilidad de la geometría, y por lo tanto sirve como banco de prueba integral de la plataforma.
